\PassOptionsToPackage{pdfpagelabels=false}{hyperref}
\documentclass[10pt,a4paper]{ctexart}

\usepackage[left=1.32cm,right=1.32cm,top=1.05cm,bottom=0.08cm]{geometry}
\usepackage{array}
\usepackage{enumitem}
\usepackage{titlesec}
\usepackage{fontawesome5}
\usepackage{xcolor}
\usepackage{hyperref}
\usepackage{bookmark}
\usepackage{eso-pic}

\definecolor{SectionBlue}{HTML}{3E68C6}
\definecolor{MutedGray}{HTML}{5B6472}
\definecolor{KeywordGray}{HTML}{1F3A52}
\definecolor{KeywordGrayZh}{HTML}{1F3A52}

% 页面设置
\pagestyle{empty}
\setlength{\parindent}{0pt}
\setlist[itemize]{leftmargin=1.15em,itemsep=0.10em,topsep=0.08em,parsep=0.01em,partopsep=0em}
\linespread{1.162}
\setmainfont{Times New Roman}
% Let ctex choose the platform-appropriate CJK fonts.

% 超链接样式
\hypersetup{
  colorlinks=true,
  urlcolor=black,
  linkcolor=black,
  pdfnewwindow=true
}

% 分节样式：标题 + 下划线
\titleformat{\section}{\heiti\zihao{3}\color{SectionBlue}}{}{0pt}{}[\vspace{0.06em}{\color{SectionBlue!40}\titlerule}]
\titlespacing*{\section}{0pt}{0.42em}{0.18em}

% 条目：左侧主标题 + 右侧角色/学位；下一行左侧描述 + 右侧时间
\newcommand{\entry}[4]{%
  \if\relax\detokenize{#3}\relax
    \textbf{{\fontsize{11.2pt}{12.2pt}\selectfont #1}}\hfill #2\enspace #4\par
  \else
    \textbf{{\fontsize{11.2pt}{12.2pt}\selectfont #1}}\hfill #2\\
    #3\hfill #4\par
  \fi
}

% 项目标题行：左侧项目名 + 中间标签 + 右侧代码/演示链接
\newcommand{\projentry}[4]{%
  \if\relax\detokenize{#2}\relax
    \makebox[0.72\textwidth][l]{\textbf{{\fontsize{11.2pt}{12.2pt}\selectfont #1}}}%
    \makebox[0.28\textwidth][r]{\codelink{#3}\if\relax\detokenize{#4}\relax\else\enspace\demolink{#4}\fi}\par
  \else
    \makebox[0.52\textwidth][l]{\textbf{{\fontsize{11.2pt}{12.2pt}\selectfont #1}}}%
    \makebox[0.20\textwidth][c]{\textcolor{MutedGray}{\fontsize{11.1pt}{11.9pt}\selectfont\textbf{#2}}}%
    \makebox[0.28\textwidth][r]{\codelink{#3}\if\relax\detokenize{#4}\relax\else\enspace\demolink{#4}\fi}\par
  \fi
}

% 仅将 GitHub 链接显示为蓝色，其他链接保持黑色
\newcommand{\ghlink}[2]{\href[pdfnewwindow=true]{#1}{\textcolor{blue}{#2}}}
\newcommand{\ghstats}[2]{\textcolor{MutedGray}{\small\faStar\ #1\enspace|\enspace\faHistory\ #2}}
\newcommand{\codelink}[1]{\href[pdfnewwindow=true]{#1}{\textcolor{blue}{\faCode\ {\fontsize{10.8pt}{11.2pt}\selectfont\textit{code}}}}}
\newcommand{\demolink}[1]{\href[pdfnewwindow=true]{#1}{\textcolor{blue}{\faVideo\ {\fontsize{10.8pt}{11.2pt}\selectfont\textit{demo}}}}}
\providecommand{\faExternalLinkAlt}{\faIcon{external-link-alt}}
\newcommand{\paperlink}[1]{\nobreak\hspace{0.18em}\href[pdfnewwindow=true]{#1}{\textcolor{MutedGray}{\fontsize{8.8pt}{9.2pt}\selectfont\faExternalLinkAlt}}}
\newcommand{\projkeyword}[1]{\textcolor{KeywordGray}{#1}}
\newcommand{\projkeywordzh}[1]{\textcolor{KeywordGrayZh}{#1}}
\newcommand{\projmetric}[1]{\textcolor{SectionBlue}{#1}}

% 匿名公开版开关：命令行定义 \PUBLICRESUME 时启用
\newif\ifpublicresume
\ifdefined\PUBLICRESUME
  \publicresumetrue
\else
  \publicresumefalse
\fi

% 教育条目：学校左对齐，专业/学历状态居中，时间右对齐
\newcommand{\eduentry}[4]{%
  \makebox[0.40\textwidth][l]{\textbf{#1}}%
  \makebox[0.25\textwidth][c]{#2}%
  \makebox[0.15\textwidth][c]{#3}%
  \makebox[0.20\textwidth][r]{#4}\par
}

\begin{document}
\AddToShipoutPictureFG*{%
  \AtPageLowerLeft{%
    \hspace*{1.32cm}\raisebox{0.25cm}[0pt][0pt]{%
      {\fontsize{8pt}{11pt}\selectfont\textcolor{MutedGray}{\begin{tabular}{@{}l@{}}
      \end{tabular}}}%
    }%
  }%
}

% ===== 顶部信息 =====
\vspace*{-1.52cm}
\begin{center}
  {\zihao{2}\bfseries \ifpublicresume Your Name (EN)\else 你的姓名\fi}\\[0.4em]
  {\fontsize{11.8pt}{12.3pt}\selectfont
  \faGlobe\ \ghlink{https://your-website.com/}{{\fontsize{10.2pt}{10.7pt}\selectfont 个人主页}}
  \enspace \enspace
  \faGithub\ \ghlink{https://github.com/your-username}{your-username}
  \enspace \ghstats{0}{0}
  \enspace \enspace
  \faEnvelope\ \href{mailto:you@example.com}{you@example.com}
  \ifpublicresume
  \else
    \enspace \enspace
    \faPhone\ +86\ 000\ 0000\ 0000
  \fi}
\end{center}
\vspace*{-0.48em}

% ===== 教育经历 =====
\section{教育背景}
\eduentry{你的大学·学院名称}{专业名称}{硕士在读}{20XX.09 -- 20XX.06}
\eduentry{你的大学·学院名称}{专业名称}{本科}{20XX.09 -- 20XX.06}

% ===== 专业技能 =====
\section{专业技能}
\begin{itemize}
  \item 理解【技术方向一】的底层原理，具备使用【工具/框架名称】等主流工具构建【应用场景】的实操经验
  \item 深入理解【技术方向二】架构与【技术方向三】技术，具备使用【框架A】与【框架B】构建复杂的【系统类型】的经验
  \item 熟悉【技术方向四】流水线，了解【技术A】与【技术B】技术，具备基于【算法名称】优化模型【能力方向】的实战经验
  \item \textbf{语言能力：}【语言考试一】\textbf{XXX} 分，【语言考试二】\textbf{XXX} 分，具备适应【语言】工作环境的能力
\end{itemize}

% ===== 项目经历 =====
\section{项目经历}
\projentry{【项目名称一】项目简短描述}{}{https://github.com/your-username/your-repo}{https://your-website.com/projects/demo}
\begin{itemize}
  \item \textbf{技术栈：}【语言】、【框架A】、【框架B】、【数据库/工具】
  \item \textbf{项目介绍：}面向【应用场景】的【系统类型】，重点解决【核心问题一】、【核心问题二】等问题。系统启动时先对【预处理步骤】进行处理，【模块A】采用【技术方案一】与【技术方案二】结合的\projkeywordzh{混合机制}。请求进入后，先对【边界条件一】、【边界条件二】等动态查询进行前置拦截，避免【不适合的场景】进入【核心流程】。随后进入\projkeywordzh{两级处理}：L1 通过【规则A】、【规则B】和【规则C】优先复用【资源类型】，L2 通过\projkeywordzh{语义处理}召回相似【资源类型】；对于存在候选但不能直接复用的情况，再由 \projkeyword{组件名称} 判断是直接复用、补充缺失信息，还是回退到完整流程。只有在【兜底条件】时，系统才进入 \projkeyword{核心算法} 式处理完成【任务类型】和【输出类型】，并将结果按规则写回【存储】。
  \item \textbf{实验评估：}在【N】条离线评测请求上，系统前置拦截【M】条【类型】查询，【指标A】覆盖率达到 \projmetric{XX.XX\%}，其中直接复用率为 \projmetric{XX.XX\%}，使端到端延迟降低了 \projmetric{XX.XX\%}，系统吞吐量提升至原来的 \projmetric{X.XX} 倍；成本方面，相对基线，平均每次请求减少 X.XX 次【调用类型】，累计净节省约【N】次调用，费用总体节省 \projmetric{XX.XX\%}。
\end{itemize}

\projentry{【项目名称二】项目简短描述}{}{https://github.com/your-username/your-repo}{https://your-website.com/projects/demo}
\begin{itemize}
  \item \textbf{技术栈：}【语言】、【框架A】、【模型A】、【模型B】、【数据库】、【工具A / 工具B】、【工具C】、【容器工具】
  \item \textbf{项目介绍：}聚焦传统【旧方案】难以稳定处理的【难点一】、【难点二】、【难点三】与【难点四】等问题，构建【新方案类型】系统。系统在上传阶段先将【输入格式A】、【输入格式B】与【输入格式C】统一转换为【中间表示】，随后使用【模型A】为每一【粒度单位】生成多向量【特征类型】，再使用【模型B】生成压缩召回向量，并将两类向量一并写入【向量数据库】建立【索引类型】。查询阶段，后端先进行【过滤步骤】，再进行\projkeywordzh{复合问题拆分}与\projkeywordzh{多轮对话上下文改写}，接着通过 \projkeyword{阶段一 + 阶段二} 精排执行两阶段检索。若高置信证据量不足，则补入 near-threshold 页面并在必要时回退到 top-1 候选页，最后通过【生成模型】生成答案，返回带【元数据A】、【元数据B】和【元数据C】溯源的结果。
  \item \textbf{实验评估：}在包含【N】份真实文档的【M】题测评中，双路检索机制展现出显著的设计优势：相比单路基线，该架构在平均检索延迟仅微增约【X】ms（至【Y】ms）的情况下，将 Hit Rate@3 由【A\%】大幅跃升至 \projmetric{XX\%}，MRR 由【B】提升至 \projmetric{X.XXXX}，并实现了 \projmetric{X.XX} 的 Faithfulness，表明系统能够以极低的时间开销换取极其稳健的定位精度。
\end{itemize}

\projentry{基于【方法A】+【方法B】的【任务名称】驱动两阶段微调}{}{https://github.com/your-username/your-repo}{}
\begin{itemize}
  \item \textbf{技术栈：}【语言】、【平台/框架A】、【技术A】、【方法A】、【方法B】
  \item \textbf{项目介绍：}针对通用模型在【问题一】、历史反馈利用不足和【问题三】的问题，围绕【具体任务】（一种【任务描述】）搭建基于 \textit{【基座模型名称】} 与【平台名称】的两阶段微调流水线。首先通过【方法A】对齐【输出格式A】和【输出格式B】双段结构化输出与基础规则推理能力，再通过【方法B】引入\projkeywordzh{奖励函数一}、\projkeywordzh{奖励函数二}、\projkeywordzh{奖励函数三}三类奖励函数，强化模型对反馈的利用和候选空间压缩能力，并构建真实对局 Benchmark 用于评估模型表现。
  \item \textbf{实验评估：}在【N】局随机基准测试、单局最多【K】次【操作】的设定下，相较基础 \textit{【基座模型】}（【A/N】通关）和单【方法B】方案（【B/N】通关，X.X/【K】步），【方法A】+【方法B】联合训练方案取得 \projmetric{X/N} 通关、获胜局平均 \projmetric{X.X/K} 步的优异成绩，表现\projkeywordzh{超过}闭源大模型\textit{【对比模型A】}（【A/N】通关，X.X/【K】步），并\projkeywordzh{逼近} \textit{【对比模型B】}（【B/N】通关，X.X/【K】步），验证了小参数模型在垂直领域推理任务上的强化学习增益。
\end{itemize}

% ===== 科研成果 =====
\section{科研成果}
\begin{itemize}
  \item \textbf{【期刊等级/会议名称】，【作者排序】：}【论文标题】
  \item \textbf{【期刊名称】（【分区/等级】），【作者排序】，【状态】：}【论文标题】
  \item \textbf{【会议名称】，【作者排序】：}【论文标题】\paperlink{https://doi.org/your-paper-doi}
  \item \textbf{一项【专利类型】}已受理（【专利申请号】）；\textbf{一项【著作权类型】}已登记（【登记号】）
\end{itemize}

% ===== 证书 =====
\section{奖项与证书}
\begin{itemize}
  \item \textbf{MOOC 认证：}自主学习【技术领域】最新技术，并在【平台A】、【平台B】、【平台C】等平台取得\ghlink{https://your-website.com/\#certificates}{多项课程证书}
  \item \textbf{竞赛：}"【竞赛主办方】"【竞赛名称】（【年份】）\textbf{【奖项等级】}（【项目名称-项目简短描述】）
  \item \textbf{奖学金：}【年份一】、【年份二】学年【奖学金等级】；本科期间获【等级一】、【等级二】奖学金各一次
\end{itemize}

\end{document}
